\section{Popolamento dell'Ontologia}
\label{sec:popolamento}

\subsection{Pipeline \texttt{populate\_wikidata.py}}

La popolazione delle istanze avviene tramite lo script \texttt{populate\_wikidata.py},
che interroga l'endpoint SPARQL pubblico di Wikidata. La pipeline \`e strutturata
in quattro fasi principali.

\subsubsection{Query SPARQL verso Wikidata}

Vengono eseguite 12 tipologie di query SPARQL, ciascuna mirata a estrarre una
specifica categoria di dati:

\begin{enumerate}[nosep]
  \item \textbf{Core} --- nomi, date di rilascio, sviluppatore e publisher (P31, P577, P178, P123);
  \item \textbf{Generi} --- mapping genere via P136;
  \item \textbf{Piattaforme} --- piattaforme via P400;
  \item \textbf{Personaggi} --- personaggi presenti nel gioco via P674;
  \item \textbf{Franchise} --- serie di appartenenza via P179;
  \item \textbf{Game Mode} --- modalit\`a di gioco via P404;
  \item \textbf{Metacritic} --- punteggio via P444 con qualificatore P447 (Metacritic);
  \item \textbf{Engine} --- motore grafico via P408;
  \item \textbf{Paese} --- paese d'origine via P495;
  \item \textbf{Sito web} --- sito ufficiale via P856;
  \item \textbf{Premi} --- premi ricevuti via P166;
  \item \textbf{Descrizioni} --- descrizioni testuali via \texttt{schema:description}.
\end{enumerate}

Per evitare i timeout dell'endpoint pubblico Wikidata (limite $\sim$60 secondi),
le query sono partizionate per \textbf{anno $\times$ mese}: 17 anni (2010--2026)
$\times$ 12 mesi = \textbf{204 iterazioni} per ciascuna tipologia di query. Tra ogni richiesta viene inserito un
\texttt{polite delay} di 1 secondo. Il tempo totale di esecuzione \`e di circa
\textbf{5 ore}.

\subsubsection{Deduplicazione}

Wikidata pu\`o restituire la stessa entit\`a con URI diversi ma stesso nome
normalizzato (es. entit\`a omonime restituite da query diverse). La pipeline
applica una normalizzazione del nome (trim degli spazi e lowercase) e rimuove
\textbf{4.420 entit\`a duplicate}, preservando per ciascun nome l'URI con il
maggior numero di triple associate.

\subsubsection{OWL-RL Reasoning}

Dopo la popolazione iniziale ($\sim$1.019.589 triple grezze), viene applicata
la chiusura deduttiva OWL 2 RL tramite \texttt{owlrl}:
\begin{itemize}[nosep]
  \item Materializzazione propriet\`a inverse: per ogni \texttt{developedBy},
    viene aggiunto \texttt{developerOf};
  \item Catene subPropertyOf: \texttt{gameName} $\rightarrow$ \texttt{name};
  \item Inferenze da classi definite (\texttt{AwardWinningGame}, \texttt{FranchiseGame}).
\end{itemize}

Il grafo cresce da $\sim$1M a $\sim$1.4M triple.

\subsubsection{Pruning}

Lo step finale di pruning (\texttt{prune\_owl.py}) rimuove triple non necessarie
a runtime:
\begin{itemize}[nosep]
  \item \texttt{owl:sameAs} riflessivi (la chiusura OWL 2 RL genera un sameAs
    di ogni entit\`a con s\'e stessa);
  \item \texttt{gameDescription}: testi lunghi non usati nelle query;
  \item \texttt{officialWebsite}: URI non necessari per la ricerca semantica.
\end{itemize}

Complessivamente vengono rimosse $\sim$253k triple.
Risultato finale: $\sim$1.173.959 triple (dataset completo 2010--2026).

\subsection{Demo Subset 2020--2026}

Per rispettare i limiti di RAM del deploy su Render (512 MB), la demo live utilizza
un sottoinsieme dell'ontologia limitato ai giochi dal 2020 in poi (\texttt{prune\_years.py}):
\begin{itemize}[nosep]
  \item $\sim$745.400 triple;
  \item $\sim$68.700 giochi;
  \item RAM a runtime: $\sim$250 MB (pyoxigraph).
\end{itemize}

\subsection{Performance: rdflib vs pyoxigraph}

La Tabella~\ref{tab:perf} confronta le performance dei due store RDF utilizzati
nel progetto. Il passaggio da \texttt{rdflib} (Python puro) a \texttt{pyoxigraph}
(binding Rust) ha prodotto un miglioramento drastico.

\begin{table}[h]
\centering
\caption{Confronto performance rdflib vs pyoxigraph}
\label{tab:perf}
\begin{tabular}{lccc}
\toprule
\textbf{Metrica}              & \textbf{rdflib} & \textbf{pyoxigraph (completo)} & \textbf{pyoxigraph (demo)} \\
\midrule
RAM a runtime                 & $\sim$1.550 MB  & $\sim$389 MB                   & $\sim$250 MB               \\
Tempo di caricamento          & $\sim$41s       & $\sim$2.2s                     & $\sim$1.5s                 \\
Query SPARQL tipica           & 1--23s          & 0.08--0.37s                    & 0.05--0.25s                \\
Speedup query                 & ---             & 5--62$\times$                  & ---                        \\
\bottomrule
\end{tabular}
\end{table}

\texttt{pyoxigraph} \citep{pyoxigraph} garantisce un caricamento quasi istantaneo
($<$3 secondi contro 41 secondi) e query 5--62 volte pi\`u veloci, rendendo
l'esperienza utente fluida anche su hardware limitato.
